% GENERATED FILE. Edit canonical lessons, then run npm run generate:books.
% canonical-source-hash: fnv1a64:101da653299e8656
% canonical-lessons: GE-C92-sauber, GE-C92-schmutzig, GE-C92-dick, GE-C92-duenn, GE-C92-leise

\chapter{Clean, Dirty, Thick, Thin, Quiet}
\label{ch:ge-clean-dirty-thick-thin-quiet}
\hlchaptermodality{92}

\begin{chapteropening}
\textbf{By the end of this chapter:} \emph{I can say clean, dirty, thick, thin and quiet.}

Complete the last lesson of chapter 92: I can say clean, dirty, thick, thin and quiet.
\end{chapteropening}

\section[sauber]{sauber — clean}

\label{lesson:GE-C92-sauber}



\begin{quote}
Before the new one: say the German for ugly, then the German for sweet.
\end{quote}

\subsection*{You'll want to know: sauber}
\textbf{sauber} — \textquotedblleft{}clean\textquotedblright{}.

\begin{cousinweb}
\textbf{sauber} $\leftarrow$ Old High German \emph{sūbar}, from Latin \emph{sōbrius}, \textquotedblleft{}sober\textquotedblright{} — an early loan. EtymDB 2.1 traces English \emph{sober} to the same Latin word.
\end{cousinweb}

\begin{grammarlens}[title={What you've built}]
One more word for this chapter.
\end{grammarlens}

\subsection*{Your turn}
\begin{itemize}
\raggedright
  \item \emph{Say it:} \emph{sauber}
  \item \emph{Say it:} \emph{sauber}, once more
  \item \emph{Say it:} say \emph{sauber}
  \item \emph{Recall:} say the German for ugly, then the German for sweet, then say \emph{sauber} again
\end{itemize}

\begin{tcolorbox}[breakable,colback=teal!4,colframe=teal!35!black,title={Before you move on}]
What does sauber mean? (\textquotedblleft{}Clean\textquotedblright{}.) Say it once more.
\end{tcolorbox}
\section[schmutzig]{schmutzig — dirty}

\label{lesson:GE-C92-schmutzig}



\begin{quote}
Before the new one: say the German for sweet, then the German for clean.
\end{quote}

\subsection*{You'll want to know: schmutzig}
\textbf{schmutzig} — \textquotedblleft{}dirty\textquotedblright{}.

\begin{cousinweb}
\textbf{schmutzig} is the word for \textquotedblleft{}dirt\textquotedblright{} + \emph{-ig}. That word $\leftarrow$ Middle High German \emph{smuz}, and EtymDB 2.1 lists it beside English \emph{smut}.
\end{cousinweb}

\begin{grammarlens}[title={What you've built}]
One more word for this chapter.
\end{grammarlens}

\subsection*{Your turn}
\begin{itemize}
\raggedright
  \item \emph{Say it:} \emph{schmutzig}
  \item \emph{Say it:} \emph{schmutzig}, once more
  \item \emph{Say it:} say \emph{schmutzig}
  \item \emph{Recall:} say the German for sweet, then the German for clean, then say \emph{schmutzig} again
\end{itemize}

\begin{tcolorbox}[breakable,colback=teal!4,colframe=teal!35!black,title={Before you move on}]
What does schmutzig mean? (\textquotedblleft{}Dirty\textquotedblright{}.) Say it once more.
\end{tcolorbox}
\section[dick]{dick — thick; fat}

\label{lesson:GE-C92-dick}



\begin{quote}
Before the new one: say the German for clean, then the German for dirty.
\end{quote}

\subsection*{You'll want to know: dick}
\textbf{dick} — \textquotedblleft{}thick; fat\textquotedblright{}.

\begin{cousinweb}
\textbf{dick} $\leftarrow$ Proto-Germanic \emph{*þekuz}, \textquotedblleft{}thick\textquotedblright{} — inherited, the same word as English \emph{thick}. (EtymDB 2.1)
\end{cousinweb}

\begin{grammarlens}[title={What you've built}]
One more word for this chapter.
\end{grammarlens}

\subsection*{Your turn}
\begin{itemize}
\raggedright
  \item \emph{Say it:} \emph{dick}
  \item \emph{Say it:} \emph{dick}, once more
  \item \emph{Say it:} say \emph{dick}
  \item \emph{Recall:} say the German for clean, then the German for dirty, then say \emph{dick} again
\end{itemize}

\begin{tcolorbox}[breakable,colback=teal!4,colframe=teal!35!black,title={Before you move on}]
What does dick mean? (\textquotedblleft{}Thick; fat\textquotedblright{}.) Say it once more.
\end{tcolorbox}
\section[dünn]{dünn — thin}

\label{lesson:GE-C92-duenn}



\begin{quote}
Before the new one: say the German for dirty, then the German for thick; fat.
\end{quote}

\subsection*{You'll want to know: dünn}
\textbf{dünn} — \textquotedblleft{}thin\textquotedblright{}.

\begin{cousinweb}
\textbf{dünn} $\leftarrow$ Old High German \emph{dunni}, \textquotedblleft{}slender, thin\textquotedblright{} $\leftarrow$ Proto-Germanic \emph{*þunnuz} — inherited, the same word as English \emph{thin}. (EtymDB 2.1)
\end{cousinweb}

\begin{grammarlens}[title={What you've built}]
One more word for this chapter.
\end{grammarlens}

\subsection*{Your turn}
\begin{itemize}
\raggedright
  \item \emph{Say it:} \emph{dünn}
  \item \emph{Say it:} \emph{dünn}, once more
  \item \emph{Say it:} say \emph{dünn}
  \item \emph{Recall:} say the German for dirty, then the German for thick; fat, then say \emph{dünn} again
\end{itemize}

\begin{tcolorbox}[breakable,colback=teal!4,colframe=teal!35!black,title={Before you move on}]
What does dünn mean? (\textquotedblleft{}Thin\textquotedblright{}.) Say it once more.
\end{tcolorbox}
\section[leise]{leise — quiet}

\label{lesson:GE-C92-leise}



\begin{quote}
Before the new one: say the German for thick; fat, then the German for thin.
\end{quote}

\subsection*{You'll want to know: leise}
\textbf{leise} — \textquotedblleft{}quiet\textquotedblright{}.

\begin{cousinweb}
\textbf{leise} $\leftarrow$ Old High German \emph{līso}, \textquotedblleft{}weak, slow, quiet\textquotedblright{} (EtymDB 2.1) — inherited.
\end{cousinweb}

\begin{grammarlens}[title={What you've built}]
One more word for this chapter.
\end{grammarlens}

\subsection*{Your turn}
\begin{itemize}
\raggedright
  \item \emph{Say it:} \emph{leise}
  \item \emph{Say it:} \emph{leise}, once more
  \item \emph{Say it:} say \emph{leise}
  \item \emph{Recall:} say the German for thick; fat, then the German for thin, then say \emph{leise} again
\end{itemize}

\begin{tcolorbox}[breakable,colback=teal!4,colframe=teal!35!black,title={Before you move on}]
What does leise mean? (\textquotedblleft{}Quiet\textquotedblright{}.) Say it once more.
\end{tcolorbox}
